\documentclass{article}
\usepackage{multicol}
\usepackage{graphicx}
\title{Page Segmentation Algorithm}
\author{Sergey Mikhno}
\date{April 2022}

\begin{document}
\maketitle

\begin{multicols}{2}
\section{Problem Description}
Sometimes we have books which are scanned or photographed. To be able to read such books on a mobile device we have to make pages readable. We have to find all the printed symbols and separate them from images. After that the page is redrawn in a larger scale and all the symbols are reflowed.

\section{Algorithm Description}

A scanned page is translated into an array of pixel intesities. For the reflowed page we create a new array. Below we have the steps necessary to reflow the page image.

\begin{enumerate}
\item Open an image file as a grayscale array.
\item Threshold the image with OTSU and BINARY\_INV.
\item Find all components containing connected non-zero pixels.
\item For every component find bounding rectangles.
\item Eliminate all rectangles contained inside others.
\item Join all intersecting rectangles.
\item Make a histogram of rectangle heights.
\item The height with the highest frequency is the text height.
\item Mark or remove all components with the hight or width \textless than $ 5 \times $ most frequent text hight.
\item For every rectangle find a neighboring one to the right, it is a nearest rectangle  intersecting or being inside of the interval of heights [y, y + height], where y is the ordinate coordinate of the component's left upper corner.
\item Create a graph, add an edge between all components and their immediate right neighbors.
\item Find connected components in that graph.
\item Join intersecting ones.
\item Those components are text lines.
\item Sort the text lines and symbols inside them.
\item Calculate the average area of the connected pixel components for every text line, let us call it $H_a$. All the intersymbol gaps bigger than $\frac{1}{2} \times H_a$ will be the interword gaps. Use the interword gaps to split the text line into words.
\item Calculate the baseline height for every text line using the histogram of lower y coordinates. The most often occurring height is the baseline. Calculate the baseline shift for every symbol.

\end{enumerate}


\end{multicols}
	

\begin{figure}
\includegraphics[width=\textwidth]{baseline.png}
\label{fig:baseline}
\caption{Baseline detection}
\end{figure}

	
\end{document}
	